\chapter{Ответы к упражнениям}\label{app:otvety}

Номер упражнения состоит из номера урока и номера внутри урока. У тех упражнений, где нужна программа, указан файл в папке \texttt{solutions}. Ваше решение может отличаться от нашего: главное, чтобы проверка проходила и шоу делало то, что требовалось.

\newcommand{\ansl}[2]{\section*{Урок #1. #2}\addcontentsline{toc}{section}{Урок #1. #2}\phantomsection\label{ans:l#1}}
\newcommand{\ans}[2]{\par\smallskip\noindent\textbf{#1.}\ #2\par}

\ansl{1}{Знакомство с окном}
\ans{1.1}{30 секунд (в управлении временем «0,0 / 30,0 с»).}
\ans{1.2}{Жёлтый (\kod{rgb('yellow')} в программе дрона 3; цвет виден и на карте, и в партитуре).}
\ans{1.3}{Клетка 49 (на вкладке «Код»: «Дрон 4 стоит на клетке 49»).}
\ans{1.4}{«Аварийная посадка: дронов на связи нет, остановлена симуляция.»}

\ansl{2}{Первый взлёт}
\ans{2.1}{Нужно заменить \kod{takeoff(100, 5)} на \kod{takeoff(60, 4)}.}
\answerfile{u02_1.py}
\ans{2.2}{Шоу длится $5 + 5 + 3 + 3 = 16$ секунд: \kod{rgb} длится ноль секунд, а ожидание три.}
\answerfile{u02_2.py}
\ans{2.3}{$5 + 8 + 3 = 16$ секунд.}
\ans{2.4}{Нужно понизить высоту до 150\cm или меньше, либо поднять лимит высоты до 180\cm и выше (в окошке на верхней панели, потом снова «Проверить»).}
\answerfile{u02_4.py}

\ansl{3}{Свет}
\ans{3.1}{Шоу длится $5 + 5 + 3 + 1 + 3 + 3 = 20$ секунд (последние 3 --- посадка).}
\answerfile{u03_1.py}
\ans{3.2}{Номера 1--7 --- это красный, зелёный, жёлтый, синий, фиолетовый, голубой, белый (имена по порядку: \kod{red}, \kod{green}, \kod{yellow}, \kod{blue}, \kod{purple}, \kod{cyan}, \kod{white}).}
\answerfile{u03_2.py}
\ans{3.3}{}
\answerfile{u03_3.py}
\ans{3.4}{Зритель увидит только зелёный. Красный включился и сразу сменился зелёным: между командами прошло ноль секунд.}
\answerfile{u03_4.py}

\ansl{4}{Высота и ожидание}
\ans{4.1}{}
\answerfile{u04_1.py}
\ans{4.2}{}
\answerfile{u04_2.py}
\ans{4.3}{}
\answerfile{u04_3.py}
\ans{4.4}{}
\answerfile{u04_4.py}

\ansl{5}{Коврик и клетки}
\ans{5.1}{$3 + 12\times 2 = 27$.}
\ans{5.2}{$40 = 12\times 3 + 4$: столбец 4, строка 3.}
\ans{5.3}{Клетка 52: столбец 4, строка 4. Центр в точке $x = 160$\cm, $y = 160$\cm.}
\ans{5.4}{Между 13 и 16: $3\times 40 = 120$\cm. Между 13 и 52 по диагонали: $\sqrt{120^2 + 120^2}\approx 170$\cm.}
\ans{5.5}{Клеток $6\times 6 = 36$. Правая верхняя: столбец 5, строка 5, номер $5 + 60 = 65$.}
\ans{5.6}{Клетка 97 в строке 8, значит над ней строка 9: $1 + 12\times 9 = 109$ (то есть 97 + 12).}

\ansl{6}{Перелёт в клетку}
\ans{6.1}{Расстояние $2\times 40 = 80$\cm, скорость $80/3\approx 27$\cm/с.}
\answerfile{u06_1.py}
\ans{6.2}{Расстояние 120\cm, за 1 секунду получается 120\cm/с. Нужно время не меньше $120/100 = 1{,}2$ секунды; берём 4 секунды.}
\answerfile{u06_2.py}
\ans{6.3}{$5 + 5 + 1 + 4 + 4 + 3 = 22$ секунды.}
\answerfile{u06_3.py}
\ans{6.4}{$170/100 = 1{,}7$ секунды. Округляем вверх до 2, но мы выбрали с запасом 4 секунды (скорость около 42\cm/с).}
\answerfile{u06_4.py}

\ansl{7}{Маршрут из нескольких клеток}
\ans{7.1}{Стороны 120, 120 и 170\cm, времена 4, 4 и 6 секунд. Шоу длится $5+5+1+4+4+6+3 = 28$ секунд.}
\answerfile{u07_1.py}
\ans{7.2}{Стороны 120, 80, 120, 80\cm; на четыре перехода 14 секунд, по 3,5. Шоу длится 28 секунд.}
\answerfile{u07_2.py}
\ans{7.3}{}
\answerfile{u07_3.py}
\ans{7.4}{}
\answerfile{u07_4.py}

\ansl{8}{Посадка в клетку}
\ans{8.1}{Шоу длится $5 + 5 + 1 + 4 + 9 = 24$ секунды (\kod{land_on} занимает 9 секунд).}
\answerfile{u08_1.py}
\ans{8.2}{Длительность увеличится на секунду: $3 + 1{,}5 + 2{,}5 + 3 = 10$ вместо 9, итого 25 секунд.}
\answerfile{u08_2.py}
\ans{8.3}{Результат тот же, 24 секунды.}
\answerfile{u08_3.py}
\ans{8.4}{}
\answerfile{u08_4.py}

\ansl{9}{Смотрим шоу: партитура}
\ans{9.1}{Команда \kod{d.takeoff(100, 5)}: взлёт идёт с 5-й по 10-ю секунду (в программе дрона это третья строка).}
\ans{9.2}{Команда \kod{d.code_s('16 52 49 13', 16)}: квадрат идёт с 11-й по 27-ю секунду.}
\ans{9.3}{На 19-й секунде: $11 + 4 + 4 = 19$.}
\ans{9.4}{30 секунд.}
\ans{9.5}{Красный (в программе дрона 2 стоит \kod{rgb('red')}).}

\ansl{10}{Два дрона}
\ans{10.1}{Пути параллельны, горизонтальное расстояние между столбцами $3\times 40 = 120$\cm, что больше 60\cm.}
\answerfile{u10_1.py}
\ans{10.2}{}
\answerfile{u10_2.py}
\ans{10.3}{Квадрат на 80\cm, висящий дрон на 150\cm: зазор 70\cm.}
\answerfile{u10_3.py}
\ans{10.4}{Дроны летят на разных высотах: 80 и 150\cm.}
\answerfile{u10_4.py}

\ansl{11}{Проверки}
\ans{11.1}{Понижены высоты до 80 и 150\cm, дрону 1 добавлено время (4 секунды вместо 1), дроны разведены по высоте.}
\answerfile{u11_1.py}
\ans{11.2}{Поднять лимит высоты до 180\cm или больше (до 200\cm) и снова нажать «Проверить». Файл менять не нужно.}
\ans{11.3}{Разбить ожидание на два шага по 45 секунд: каждый короче 60.}
\answerfile{u11_3.py}
\ans{11.4}{Клетка 8 лежит в столбце 8, а четверть 1 --- это столбцы 0--5. Заменяем на клетку 5.}
\answerfile{u11_4.py}

\ansl{12}{Очередь и ожидание}
\ans{12.1}{Дрон 2 ждёт 9 секунд (8 секунд полёта дрона 1 и запас).}
\answerfile{u12_1.py}
\ans{12.2}{}
\answerfile{u12_2.py}
\ans{12.3}{}
\answerfile{u12_3.py}
\ans{12.4}{«Позже» записывается ожиданием: перед движением \kod{set_z(100, 3)} на той же высоте.}
\answerfile{u12_4.py}

\ansl{13}{Повторы: циклы}
\ans{13.1}{}
\answerfile{u13_1.py}
\ans{13.2}{}
\answerfile{u13_2.py}
\ans{13.3}{Дрон \kod{n} ждёт \kod{2 * n} секунд до своей вспышки и \kod{2 * (2 - n)} после неё, поэтому длина у всех одинаковая.}
\answerfile{u13_3.py}
\ans{13.4}{}
\answerfile{u13_4.py}

\ansl{14}{Другая четверть}
\ans{14.1}{Клетки 13, 16, 52, 49 плюс 72: 85, 88, 124, 121.}
\answerfile{u14_1.py}
\ans{14.2}{Клетки 13, 16, 52, 49 плюс 78: 91, 94, 130, 127.}
\answerfile{u14_2.py}
\ans{14.3}{$14 + 72 = 86$; $55 + 78 = 133$.}
\ans{14.4}{$100 = 12\times 8 + 4$: столбец 4, строка 8. Столбцы 0--5 и строки 6--11: это четверть 3. В четверти 1 то же место имеет номер $100 - 72 = 28$.}

\ansl{15}{Круги и подъём на ходу}
\ans{15.1}{Меняем \kod{'cw'} на \kod{'ccw'}.}
\answerfile{u15_1.py}
\ans{15.2}{Дрон оказывается в клетке 38 (на противоположной стороне круга: центр клетки 26, клетка 14 снизу, клетка 38 сверху).}
\answerfile{u15_2.py}
\ans{15.3}{}
\answerfile{u15_3.py}
\ans{15.4}{Дроны стоят на противоположных концах диаметра длиной 80\cm и крутятся с одной скоростью в одну сторону, расстояние между ними всё время 80\cm, больше порога 60\cm.}
\answerfile{u15_4.py}

\ansl{16}{Итоговые проекты}
\ans{Проект 1}{}
\answerfile{t06_follow.py}
\ans{Проект 2}{}
\answerfile{t09_ladder.py}
\ans{Проект 3}{}
\answerfile{t07_corners.py}
\ans{Проект 4}{}
\answerfile{t15_star5.py}
\ans{Проект 5}{}
\answerfile{t16_rows.py}
\ans{Проект 6}{}
\answerfile{t17_hex7.py}
\ans{Проект 7}{}
\answerfile{t18_eight.py}
